% !TEX program = lualatex

\documentclass{chamranintern}

%% ------------------------------ اطلاعات روی جلد (این بخش را ویرایش کنید) ----
\ReportTitle{عنوان گزارش کارآموزی}
\StudentName{<اسم دانشجو>}
\StudentID{شماره دانشجو}
\Supervisor{استاد راهنما(کارآموزی)}
\Organization{محل کارآموزی}
\InternPeriod{تاریخ شروع و پایان}
\IndustrySupervisor{سرپرست صنعتی}   % اختیاری
\Department{گروه مهندسی کامپیوتر}        % اختیاری
\ReportDate{اهواز، ۱۴۰۵}                 % اختیاری
\Keywords{اینترنت اشیا، [حوزه صنعتی]، [نوع مسأله]، \ldots}

\begin{document}

%% ----------------------------------------------------------------- جلد
\maketitle

%% ----------------------------------------------------------------- چکیده
\begin{abstract}
  راهنما: چکیده باید ۲۰۰ تا ۳۰۰ کلمه باشد و موارد زیر را پوشش دهد:
  \begin{itemize}
    \item معرفی کوتاه سازمان/صنعت؛
    \item مسأله‌ی شناسایی‌شده (به‌صورت خیلی مختصر)؛
    \item راه‌حل پیشنهادی مبتنی بر \lat{IoT}؛
    \item مهم‌ترین چالش‌ها و محدودیت‌ها؛
    \item دستاورد اصلی کارآموزی.
  \end{itemize}
\end{abstract}

%% ------------------------------------------------------------- فهرست‌ها
\tableofcontents
% اگر در گزارش شکل یا جدول دارید، دو خط زیر را فعال کنید:
% \listoffigures
% \listoftables

%% ===================================================================
\chapter{مقدمه}
\section{مقدمه}
توضیح کلی درباره‌ی اهمیت اینترنت اشیا در صنعت/سازمان موردنظر و ضرورت انجام این کارآموزی.

\section{بیان ضرورت و اهمیت موضوع}
چرا این مسأله مهم است؟ تأثیر آن بر بهره‌وری، هزینه، کیفیت، ایمنی و \ldots

\section{اهداف کارآموزی}
\begin{itemize}
  \item هدف اصلی
  \item اهداف فرعی (شناسایی مسأله، ارائه راه‌حل، بررسی چالش‌ها و \ldots)
\end{itemize}

\section{ساختار گزارش}
معرفی کوتاه فصل‌های بعدی.

%% ===================================================================
\chapter{معرفی سازمان/صنعت محل کارآموزی}
\section{معرفی کلی سازمان}
\begin{itemize}
  \item نام، حوزه‌ی فعالیت، اندازه، تاریخچه‌ی مختصر
  \item ساختار سازمانی مرتبط با حوزه‌ی کارآموزی
\end{itemize}

\section{فرآیندها و سیستم‌های فعلی}
\begin{itemize}
  \item توصیف فرآیندهای کلیدی مرتبط با مسأله
  \item فناوری‌های موجود (اتوماسیون، نرم‌افزارها، سنسورها و \ldots)
\end{itemize}

\section{وضعیت موجود (از منظر فناوری‌های مرتبط با مسأله)}
\begin{itemize}
  \item زیرساخت‌ها و فناوری‌های موجود
\end{itemize}

%% ===================================================================
\chapter{مرور پیشینه و فناوری‌های مرتبط با پژوهش}
\section{کاربردهای \lat{IoT} در صنعت/حوزه‌ی مشابه}
\begin{itemize}
  \item مرور مقالات، پروژه‌ها و نمونه‌های موفق \cite{atzori2010iot,gubbi2013iot}
  \item جدول مقایسه‌ای راه‌حل‌های احتمالی
\end{itemize}

\section{جمع‌بندی و شکاف پژوهشی/کاربردی}
\begin{itemize}
  \item چه چیزی در سازمان محل کارآموزی وجود ندارد یا ضعیف است؟
\end{itemize}

%% ===================================================================
\chapter{شرح انجام کارآموزی}
\section{رویکرد انجام کار}
\begin{itemize}
  \item نحوه‌ی انجام کار (کارهای مطالعاتی، میدانی، عملیاتی، شبیه‌سازی و \ldots)
  \item مراحل انجام کار (نمودار بلوکی - جریانی)
\end{itemize}

\section{ابزارها و روش‌های جمع‌آوری داده و اطلاعات}

%% ===================================================================
\chapter{شناسایی و تعریف دقیق مسأله}
\section{توصیف وضعیت موجود}
\begin{itemize}
  \item شرح فرآیند فعلی به‌صورت دقیق
\end{itemize}

\section{شناسایی نقاط ضعف و چالش‌ها}
\begin{itemize}
  \item مشکلات عملیاتی، هزینه‌ای، ایمنی، کیفیت، زمان و \ldots
  \item شواهد و داده‌های پشتیبان
\end{itemize}

\section{تعریف دقیق مسأله}
\begin{itemize}
  \item \textbf{بیان مسأله} در یک پاراگراف شفاف
  \item دامنه‌ی مسأله: آنچه داخل و خارج از محدوده است.
  \item ذی‌نفعان مسأله (داخل و بیرون صنعت چه کسانی از حل این مشکل سود می‌برند)
  \item معیارهای موفقیت (\lat{KPIs}): چه معیارهای قابل اندازه‌گیری برای اعتبارسنجی حل مسأله وجود دارد؟
\end{itemize}

%% ===================================================================
\chapter{راه‌حل پیشنهادی مبتنی بر اینترنت اشیا}
\section{معماری کلی راه‌حل}
\begin{itemize}
  \item نمودار معماری (لایه‌ها، اجزا، ارتباطات)
  \item توضیح جریان داده از دریافت تا تصمیم‌گیری
\end{itemize}

\section{فناوری‌های احتمالی مورد استفاده راه‌حل}
\begin{table}[h]
  \centering
  \caption{فناوری‌های پیشنهادی برای هر بخش از راه‌حل}
  \renewcommand{\arraystretch}{1.6}
  \begin{tabular}{|p{3.6cm}|p{4.4cm}|p{6.2cm}|}
    \hline
    \textbf{بخش} & \textbf{فناوری پیشنهادی} & \textbf{دلیل انتخاب} \\ \hline
    سنسورها و عملگرها &  &  \\ \hline
    شبکه ارتباطی &  &  \\ \hline
    بستر پردازش (\lat{Cloud/Edge}) &  &  \\ \hline
    پلتفرم \lat{IoT} &  &  \\ \hline
    لایه تحلیل و هوش &  &  \\ \hline
    رابط کاربری &  &  \\ \hline
  \end{tabular}
\end{table}

\section{الزامات پیاده‌سازی}
\begin{itemize}
  \item سخت‌افزاری
  \item نرم‌افزاری
  \item ارتباطی
  \item نیروی انسانی و آموزش
\end{itemize}

%% ===================================================================
\chapter{چالش‌ها و محدودیت‌های اجرا}
\section{چالش‌های فنی}
\begin{itemize}
  \item محدودیت پهنای باند، توان مصرفی، تداخل، مقیاس‌پذیری، امنیت و \ldots
\end{itemize}

\section{چالش‌های سازمانی و انسانی}
\begin{itemize}
  \item مقاومت در برابر تغییر، کمبود مهارت، محدودیت‌های سازمانی و قانونی، محدودیت‌های امنیت و حریم خصوصی
\end{itemize}

\section{چالش‌های مالی}
\begin{itemize}
  \item هزینه‌های سرمایه‌ای و جاری، توجیه اقتصادی
\end{itemize}

\section{محدودیت‌های عملیاتی}
\begin{itemize}
  \item محدودیت زمانی، دسترسی به داده، عدم امکان پیاده‌سازی کامل
\end{itemize}

\section{راهکارهای کاهش چالش‌ها (پیشنهادی)}

%% ===================================================================
\chapter{جمع‌بندی و پیشنهادها}

%% ------------------------------------------------------------ منابع
\bibliographystyle{unsrt}
\bibliography{references}

\end{document}
